\documentclass[10pt]{article}
\usepackage[letterpaper,margin=0.62in]{geometry}
\usepackage{graphicx,xcolor,enumitem,hyperref,multicol,tabularx}
\usepackage[T1]{fontenc}
\usepackage{lmodern}
\pagestyle{empty}
\setlength{\parindent}{0pt}
\setlength{\parskip}{4pt}
\setlist[itemize]{leftmargin=1.15em,itemsep=1.5pt,topsep=2pt}
\definecolor{askred}{HTML}{7B2933}
\definecolor{askink}{HTML}{21191B}
\hypersetup{colorlinks=true,urlcolor=askred,linkcolor=askred}
\newcommand{\head}[1]{\vspace{4pt}{\large\bfseries\color{askred}#1}\par\vspace{1pt}}

\begin{document}

\begin{center}
  \includegraphics[width=\textwidth,trim=0 280 0 280,clip]{../assets/ask-2027-hero.jpg}
  \vspace{5pt}

  {\fontsize{24}{27}\selectfont\bfseries\color{askink} ASK 2027}\par
  {\fontsize{15}{18}\selectfont\bfseries Workshop on Agentic Methods for Sparse and Irregular HPC Kernels}\par
  \vspace{3pt}
  {\bfseries Wednesday morning, June 2, 2027 \textbullet{} IPDPS 2027 \textbullet{} Seattle, Washington}\par
  \href{https://ypsoh.github.io/ask-ipdps-2027/}{https://ypsoh.github.io/ask-ipdps-2027/}
\end{center}

\small
\head{Call for Papers}
ASK 2027 invites original, unpublished work on agentic methods for high-performance computing
kernels, with particular emphasis on sparse and irregular workloads. The workshop brings together
researchers in sparse and combinatorial computation, compilers and performance tools, and agentic
AI systems. Negative results and failure analyses are explicitly welcome.

\begin{multicols}{2}
\raggedcolumns
\head{Topics include}
\begin{itemize}
  \item Compile--run--profile--revise optimization loops
  \item Counter- and roofline-guided agentic search
  \item Sparse format and algorithm selection
  \item Input-dependent kernel specialization
  \item Sparse tensor kernels and decompositions
  \item Multi-agent generation, verification, tuning, and porting
  \item Compiler, profiler, performance-model, and sparse-DSL tool use
  \item Correctness and numerical verification
  \item Performance portability across CPUs, GPUs, and accelerators
  \item Benchmarks and evaluation for HPC coding agents
  \item Cost, energy, and human-in-the-loop tradeoffs
\end{itemize}
Dense, stencil, and FFT kernels are in scope when the primary contribution is the agentic method.

\head{Submission tracks}
\textbf{Full papers:} up to 10 pages for complete research results and substantial systems,
methods, benchmarks, or empirical studies.

\textbf{Short papers:} up to 4 pages for work in progress, position papers, experience reports,
negative results, and emerging ideas. Accepted short papers will also be presented at the
cross-workshop poster session.

\head{Formatting}
Use the official \href{https://www.ieee.org/conferences/publishing/templates.html}{IEEE conference
proceedings template}: US Letter (8.5 $\times$ 11 inches), 10-point type, single-spaced,
double-column. Page limits include the body, figures, and tables; references are excluded.
Appendices must fit within the applicable limit. Submit one PDF through Linklings.

ASK uses \textbf{single-anonymous review}; include author names, affiliations, and acknowledgments.
Submissions must not be published or simultaneously under review at another archival venue.

\head{Review and publication}
Every paper will receive at least three reviews. Accepted full and short papers will appear in the
IPDPS Workshops proceedings (IPDPSW), submitted for inclusion in IEEE Xplore. At least one author
of each accepted paper must register and present in person.

Authors should report model and tool versions, sampling parameters, optimization budgets,
hardware/software environments, correctness criteria, and the number of trials. Code, matrices,
prompts, agent trajectories, and other artifacts are strongly encouraged.

\head{Important dates}
\begin{tabular}{@{}>{\bfseries}l l@{}}
Submission & Jan. 29, 2027, 11:59 p.m. AoE \\
Notification & Feb. 22, 2027 \\
Camera-ready & Mar. 6, 2027 \\
Poster session & June 1, 2027, 6:00--8:00 p.m. \\
Workshop & June 2, 2027, morning \\
\end{tabular}

\head{Organizers and contact}
Ramakrishnan Kannan (ORNL), Yongseok Soh (ORNL), and Ariful Azad (Texas A\&M University).\\
Questions: \href{mailto:sohyp@ornl.gov}{sohyp@ornl.gov}. The workshop-specific Linklings URL will
be posted on the workshop website.
\end{multicols}

\end{document}
